\documentclass[11pt]{article}

\usepackage[margin=1in]{geometry}
\usepackage{amsmath, amssymb}
\usepackage{xcolor}
\usepackage{listings}
\usepackage{hyperref}
\usepackage{array}
\usepackage{booktabs}

\hypersetup{
    colorlinks=true,
    linkcolor=blue!60!black,
    urlcolor=blue!60!black,
    citecolor=blue!60!black
}

\lstdefinestyle{bugpython}{
  language=Python,
  basicstyle=\ttfamily\small,
  keywordstyle=\color{blue!60!black},
  stringstyle=\color{green!40!black},
  commentstyle=\color{gray!70!black},
  showstringspaces=false,
  breaklines=true,
  frame=single,
  rulecolor=\color{gray!30},
  columns=fullflexible,
  keepspaces=true
}

\title{SymPy Bug Report}
\author{}
\date{}

\begin{document}

\maketitle

\section{Bug 8: K-Weighted Alternating Binomial Sum Misses the \texorpdfstring{\(n=1\)}{n=1} Case}

\subsection{Minimal Reproducer}

\medskip

\begin{lstlisting}[style=bugpython]
from sympy import binomial, summation, symbols

n = symbols("n", integer=True, nonnegative=True)
k = symbols("k", integer=True)

closed = summation((-1)**k*k*binomial(n, k), (k, 0, n))
direct_n1 = sum(((-1)**j)*j*binomial(1, j) for j in range(0, 2))

print("closed form =", closed)
print("closed at n=1 =", closed.subs(n, 1))
print("direct finite sum at n=1 =", direct_n1)
print("difference =", closed.subs(n, 1) - direct_n1)
\end{lstlisting}

\subsection{Observed Behavior}

\medskip

\begin{lstlisting}[style=bugpython]
closed form = 0
closed at n=1 = 0
direct finite sum at n=1 = -1
difference = 1
\end{lstlisting}

SymPy returns the unconditional closed form \(0\).  Substituting \(n=1\) into
that result gives \(0\), while direct evaluation of the finite sum at \(n=1\)
gives \(-1\).  The returned scalar therefore loses a parameter value allowed
by the declared assumption \(n \in \mathbb{Z}_{\ge 0}\).

\subsection{Expected Behavior}

For a nonnegative integer symbol \(n\), the result should preserve the
exceptional \(n=1\) value.  A correct result for the representative sum is
\[
  \operatorname{Piecewise}\bigl((-1,\operatorname{Eq}(n,1)),(0,\operatorname{Ne}(n,1))\bigr),
  \qquad n\in\mathbb{Z}_{\ge 0},
\]
or any equivalent expression that evaluates to \(-1\) at \(n=1\) and \(0\) at
all other nonnegative integer values of \(n\).

\subsection{Mathematical Explanation}

Differentiate the binomial identity
\[
  (1+t)^n=\sum_{k=0}^{n}\binom{n}{k}t^k
\]
with respect to \(t\).  This gives
\[
  n(1+t)^{n-1}=\sum_{k=0}^{n} k\binom{n}{k}t^{k-1}.
\]
Multiplying by \(t\) and setting \(t=-1\) yields
\[
  \sum_{k=0}^{n}(-1)^k k\binom{n}{k}
  = -n(1-1)^{n-1}.
\]
For \(n>1\), this is \(0\).  For \(n=1\), however, the differentiated
identity must be evaluated as the original finite sum:
\[
  \sum_{k=0}^{1}(-1)^k k\binom{1}{k}
  = 0\binom{1}{0} - 1\binom{1}{1}
  = -1.
\]
The case \(n=0\) is also \(0\), since the only summand contains the factor
\(k=0\).  Thus the scalar \(0\) is not an equivalent simplification over the
declared nonnegative-integer domain; it has the wrong value at \(n=1\).

\subsection{Source-Level Diagnosis}

The diagnosis locates the failure in SymPy's concrete summation path.  The
public call constructs a \texttt{Sum} and calls \texttt{doit}.  The relevant
source files are:

\begin{center}
\begin{tabular}{@{}l@{}}
\nolinkurl{sympy/concrete/summations.py}\\
\nolinkurl{sympy/concrete/gosper.py}
\end{tabular}
\end{center}

Since the range length is symbolic, evaluation reaches \texttt{eval\_sum} and
then \texttt{eval\_sum\_\allowbreak symbolic}.  The traced call path reaches
the Gosper branch around \nolinkurl{sympy/concrete/summations.py} lines
1236--1256, which calls \texttt{gosper\_sum}.

For the summand \((-1)^k k\binom{n}{k}\), the traced Gosper certificate is
\[
  g = -\frac{k-1}{n-1}.
\]
This certificate has a pole at the parameter value \(n=1\).  The diagnosis
identifies the immediate faulty endpoint calculation in
\nolinkurl{sympy/concrete/gosper.py}:

\begin{lstlisting}[style=bugpython]
result = (f*(g + 1)).subs(k, b) - (f*g).subs(k, a)
\end{lstlisting}

\texttt{gosper\_sum} substitutes the symbolic endpoints and simplifies the
generic telescoping result to the scalar \(0\), without recording that the
certificate is invalid at \(n=1\).  The caller has logic to wrap some results
in \texttt{Piecewise} when visible denominators remain in \texttt{Mul} or
\texttt{Add} results, but here the endpoint expression has already simplified
to \(0\).  As a result, the excluded parameter value is no longer visible to
\texttt{eval\_sum\_\allowbreak symbolic}, and the generic value becomes the
final answer.

The suggested fix direction is to track parameter denominators introduced by
\texttt{gosper\_term} or \texttt{gosper\_sum}, and either emit a conditional
result or fall back to direct evaluation at excluded parameter values.  For
this example, the final result needs an \(n=1\) branch.

\subsection{Additional Failing Instances}

The independent verification checks the same endpoint failure for
\[
  \sum_{k=0}^{n} (-1)^k k^p \binom{n}{k},
  \qquad p=1,2,3,4,5,6.
\]
For each listed power, SymPy returns a closed form whose value at \(n=1\) is
\(0\).  The direct finite sum at \(n=1\) is the same in all cases with
\(p\ge 1\), because the \(k=0\) term is \(0^p\binom{1}{0}=0\) and the \(k=1\)
term is \((-1)1^p\binom{1}{1}=-1\).

\begin{center}
\small
\renewcommand{\arraystretch}{1.3}
\begin{tabular}{@{}>{\raggedright\arraybackslash}p{0.44\linewidth}>{\raggedright\arraybackslash}p{0.23\linewidth}>{\raggedright\arraybackslash}p{0.23\linewidth}@{}}
\toprule
\textbf{Input / Expression} & \textbf{SymPy behavior} & \textbf{Expected behavior} \\
\midrule
\(\sum_{k=0}^{n}(-1)^k k\binom{n}{k}\), evaluated at \(n=1\) & Closed form gives \(0\) & Direct sum gives \(-1\) \\
\(\sum_{k=0}^{n}(-1)^k k^2\binom{n}{k}\), evaluated at \(n=1\) & Closed form gives \(0\) & Direct sum gives \(-1\) \\
\(\sum_{k=0}^{n}(-1)^k k^3\binom{n}{k}\), evaluated at \(n=1\) & Closed form gives \(0\) & Direct sum gives \(-1\) \\
\(\sum_{k=0}^{n}(-1)^k k^4\binom{n}{k}\), evaluated at \(n=1\) & Closed form gives \(0\) & Direct sum gives \(-1\) \\
\(\sum_{k=0}^{n}(-1)^k k^5\binom{n}{k}\), evaluated at \(n=1\) & Closed form gives \(0\) & Direct sum gives \(-1\) \\
\(\sum_{k=0}^{n}(-1)^k k^6\binom{n}{k}\), evaluated at \(n=1\) & Closed form gives \(0\) & Direct sum gives \(-1\) \\
\bottomrule
\end{tabular}
\end{center}

\subsection{Related Issues Assessment}

The bug-finding harness performed a best-effort deduplication check.  Besides
general documentation and background for Gosper summation, the open issue
\href{https://github.com/sympy/sympy/issues/29731}{sympy/sympy\#29731}
(``A summation that evaluates to 0 even though it can be nonzero depending on
the variable'', created 2026-05-02) is a near-exact public report of this same
behavior: its reproducer returns the scalar \(0\) for an \((-1)\)-weighted
\(i^k\) binomial sum that is nonzero at the singular parameter value, with the
same \texttt{gosper\_sum} lost-endpoint-exception root cause.  That issue
predates this bundle, so the deduplication verdict is a new instance of an
already-reported failure family: symbolic Gosper summation can require
parameter conditions, and this candidate contributes the explicit
\((-1)^k k\binom{n}{k}\) instantiation, the missing \(n=1\) branch, and the
\(k^p\) generalization, all of which remain individually actionable as a
cross-link to \#29731 and as a regression test.

\subsection{Proposed SymPy Regression Test}

\medskip

\begin{lstlisting}[style=bugpython]
import pytest
from sympy import binomial, summation, symbols

n = symbols("n", integer=True, nonnegative=True)
k = symbols("k", integer=True)
CASES = [1, 2, 3, 4, 5, 6]

@pytest.mark.parametrize("power", CASES)
def test_weighted_alternating_binomial_sum_preserves_n1_exception(power):
    closed = summation((-1)**k*k**power*binomial(n, k), (k, 0, n))
    direct = sum(((-1)**j)*(j**power)*binomial(1, j) for j in range(2))
    assert closed.subs(n, 1) == direct
\end{lstlisting}

\end{document}
